% Subsection: The Power Set as the Ambient Universe  (notes stub; content authored later)
\subsection{The Power Set as the Ambient Universe}

\begin{remark*}[Exposition]
Once a base set \(X\) is fixed, the power set \(\mathcal{P}(X)\) becomes the
ambient universe for this chapter.  Every family of subsets of \(X\) lives
inside \(\mathcal{P}(X)\).  The full power set is usually too large to be the
structure of interest, but it supplies the surrounding arena in which smaller
families can be compared, enlarged, intersected, and closed under operations.
\end{remark*}

\begin{remark*}[Ambient Set Convention]
Complements are always relative to the fixed base set.  If
\(\mathcal{F}\subseteq\mathcal{P}(X)\) and \(A\in\mathcal{F}\), then the
notation \(A^c\) means
\[
  A^c = X\setminus A.
\]
Changing the ambient set changes the complement operation.  Thus the same
subset \(A\) may have different complements when it is considered inside
different base sets.
\end{remark*}

\begin{remark*}[Structural View]
A subset family is not only a list of subsets.  It is a candidate domain for
set operations.  One can ask whether \(A\cup B\), \(A\cap B\), \(A^c\), or
\(\bigcup_{n=1}^{\infty}A_n\) remains in the same family after the operation
is performed.  Those closure questions are the mechanism by which raw subset
families become mathematical structures.
\end{remark*}

\begin{remark*}[Boundary Cases]
For every set \(X\), the two extreme families
\[
  \varnothing
  \qquad\text{and}\qquad
  \mathcal{P}(X)
\]
are both families of subsets of \(X\).  They play different roles.  The empty
family contains no admitted subsets at all.  The power set admits every subset
of \(X\), so every set operation whose output is a subset of \(X\) stays inside
\(\mathcal{P}(X)\).
\end{remark*}
